% ---
% title: The binomial tree
% tags: finance, combinations, counting
% ---

\begin{problem}
	\begin{context}
		In the \emph{binomial (Cox--Ross--Rubinstein) model} of option
		pricing a stock price is modeled in discrete steps: in each step it is
		multiplied either by an up-factor $u>1$ or by a down-factor $d<1$. An
		option's value depends on the final price, so we need to know how many
		paths lead to each final price.
	\end{context}
	A stock moves up (factor $u$) or down (factor $d$) every day. After $n=10$
	days the price is $S_0u^kd^{10-k}$, where $k$ is the number of up-days.
	(a) How many price paths are there? (b) How many paths end with $k=6$
	up-days? (c) How many different final prices are there?
\end{problem}

\begin{slides}
	\small
	\begin{columns}[T]
		\begin{column}{.6\textwidth}
			\exstatement
		\end{column}
		\begin{column}{.38\textwidth}
			\centering
			\begin{tikzpicture}[scale=.75,every node/.style={font=\scriptsize}]
				\foreach \t in {0,...,3}
				\foreach \k in {0,...,\t} {
						\pgfmathtruncatemacro{\c}{factorial(\t)/(factorial(\k)*factorial(\t-\k))}
						\node[circle,draw=navy,fill=navy!8,inner sep=1.5pt] (n\t\k)
						at (1.4*\t,\k-\t/2) {\c};
					}
				\foreach \t in {0,...,2} {
						\pgfmathtruncatemacro{\tn}{\t+1}
						\foreach \k in {0,...,\t} {
								\pgfmathtruncatemacro{\kn}{\k+1}
								\draw[->,gray] (n\t\k)--(n\tn\kn);
								\draw[->,gray] (n\t\k)--(n\tn\k);
							}
					}
			\end{tikzpicture}\\
			\scriptsize paths reaching each node; each node $=$ sum of the two
			nodes feeding it (Pascal's rule)
		\end{column}
	\end{columns}
	\vspace{-1.5mm}
	\uncover<2->{\sol
			\begin{enumerate}\setlength{\itemsep}{0pt}
				\item[(a)] 2 choices on each of 10 days: $2^{10}=1024$.
				\item[(b)] A path is a word of 6 U's and 4 D's; choose the 6 up-days:
				      $\binom{10}{6}=210$.
				\item[(c)] The order of the factors does not change the product, so the
				      final price depends only on $k\in\{0,\dots,10\}$: $11$ prices.
			\end{enumerate}
			Check: grouping the paths by $k$ gives $\sum_{k=0}^{10}\binom{10}{k}=2^{10}$.
			This is the combinatorics behind the Cox--Ross--Rubinstein model.}
\end{slides}

\begin{solution}
	\idea A path is a word of 10 letters U/D. The final price depends only on
	how many U's the word contains, not on where they are.

	\step{(a)} 2 choices on each of the 10 days: $2^{10}=1024$.

	\step{(b)} A path with 6 up-days is a word with 6 U's and 4 D's; choose the
	positions of the U's: $\binom{10}{6}=210$.

	\step{(c)} The product $S_0u^kd^{10-k}$ does not depend on the order of the
	factors, so the final price is determined by $k\in\{0,1,\dots,10\}$: 11
	possible prices (they are different because $u>d$).

	\begin{center}
		\begin{tikzpicture}[scale=.85,every node/.style={font=\scriptsize}]
			\foreach \t in {0,...,4}
			\foreach \k in {0,...,\t} {
					\pgfmathtruncatemacro{\c}{factorial(\t)/(factorial(\k)*factorial(\t-\k))}
					\node[circle,draw=navy,fill=navy!8,inner sep=1.5pt,minimum size=13pt]
					(n\t\k) at (1.5*\t,\k-\t/2) {\c};
				}
			\foreach \t in {0,...,3} {
					\pgfmathtruncatemacro{\tn}{\t+1}
					\foreach \k in {0,...,\t} {
							\pgfmathtruncatemacro{\kn}{\k+1}
							\draw[->,gray] (n\t\k)--(n\tn\kn);
							\draw[->,gray] (n\t\k)--(n\tn\k);
						}
				}
			\node[right,align=left,font=\small] at (6.6,0)
			{number of paths reaching each node;\\each node is the sum of the\\two
				feeding it (Pascal's rule)};
		\end{tikzpicture}
	\end{center}

	\begin{answer}
		(a) $2^{10}=1024$ \quad (b) $\binom{10}{6}=210$ \quad (c) $11$
	\end{answer}
	\begin{remark}
		Grouping the 1024 paths by their final node gives
		\[ \sum_{k=0}^{10}\binom{10}{k}=2^{10}. \]
		In the option-pricing model each path with $k$ up-moves gets
		weight $p^k(1-p)^{10-k}$, and $\binom{10}{k}$ tells how many such paths
		there are --- this is the binomial distribution we will meet later.
		In the figure every node is the sum of the two nodes feeding it: this is
		Pascal's rule of Exercise~\exref{pascals-rule}, now read as ``paths into a node''.
	\end{remark}
\end{solution}
